% Clase 2 — Lógica de Primer Orden
% Lectura: AIMA Cap. 8
\section{Lógica de Primer Orden (LPO)}

\subsection{Sintaxis}

Elementos de LPO:
\begin{itemize}
    \item \textbf{Constantes}: $Juan$, $UNIANDES$, \ldots
    \item \textbf{Variables}: $x, y, z, \ldots$
    \item \textbf{Predicados}: $Estudiante(x)$, $Mayor(x, y)$
    \item \textbf{Funciones}: $Padre(x)$, $Suma(x, y)$
    \item \textbf{Cuantificadores}: $\forall$ (universal), $\exists$ (existencial)
    \item \textbf{Conectivos}: $\neg, \land, \lor, \Rightarrow, \Leftrightarrow$
\end{itemize}

\subsection{Semántica}

Un \textbf{modelo} en LPO especifica:
\begin{itemize}
    \item El dominio de objetos
    \item La interpretación de constantes, predicados y funciones
\end{itemize}

\subsection{Cuantificadores}

$$\forall x\ P(x) \equiv \neg \exists x\ \neg P(x)$$
$$\exists x\ P(x) \equiv \neg \forall x\ \neg P(x)$$

\begin{remark}
$\forall x\ (P(x) \Rightarrow Q(x))$ es diferente de $\forall x\ (P(x) \land Q(x))$.
El cuantificador universal normalmente usa $\Rightarrow$; el existencial usa $\land$.
\end{remark}

\subsection{Ejemplo: Modelo en LPO}

\begin{example}[Modelo simple]
    Sea el dominio $D = \{\text{Juan}, \text{Maria}, \text{Pedro}\}$ y el predicado binario $Padre(x,y)$ ("$x$ es
    padre de $y$"). La interpretación de la Figura siguiente satisface $Padre(\text{Juan},\text{Maria})$ y
    $Padre(\text{Juan},\text{Pedro})$, y ninguna otra pareja.
\end{example}

\begin{figure}[H]
\centering
\begin{tikzpicture}[
  obj/.style={draw=black, thick, circle, minimum size=11mm, font=\small}
]
\node[obj] (juan)  at (0,0)   {Juan};
\node[obj] (maria) at (3.2,1) {Maria};
\node[obj] (pedro) at (3.2,-1){Pedro};

\draw[-{Latex},thick] (juan) -- node[above, font=\scriptsize]{$Padre$} (maria);
\draw[-{Latex},thick] (juan) -- node[below, font=\scriptsize]{$Padre$} (pedro);
\end{tikzpicture}
\caption{Modelo con dominio $\{\text{Juan},\text{Maria},\text{Pedro}\}$: cada flecha es una pareja $(x,y)$ para la
cual $Padre(x,y)$ es verdadero en este modelo. Cambiar las flechas cambia el modelo sin cambiar la sintaxis de la
oración.}
\end{figure}

\begin{exercise}
    Usando los predicados $Estudiante(x)$, $Curso(y)$, $Aprueba(x,y)$ y $Certificado(x,y)$, traduzca a LPO:
    \begin{enumerate}
        \item Todo estudiante que aprueba un curso recibe un certificado de ese curso.
        \item Existe al menos un estudiante que no ha aprobado ningún curso.
        \item Ningún estudiante aprueba todos los cursos.
    \end{enumerate}
    Para cada oración, indique si el cuantificador principal es $\forall$ o $\exists$ y justifique si el conectivo
    correcto dentro de su alcance es $\Rightarrow$ o $\land$.
\end{exercise}
